\documentclass[11pt]{article}
\usepackage[margin=1in]{geometry}
\usepackage[T1]{fontenc}
\usepackage{lmodern}
\usepackage[hidelinks]{hyperref}
\usepackage{booktabs}
\usepackage{microtype}

\title{Tactix Documentation Maintainability}
\author{Tactix Development Project}
\date{August 8, 2026}

\begin{document}
\maketitle

\section*{Executive summary}
The project's living documentation has been reorganized around one rule: each
fact has one authoritative home. The plan now contains stable commitments,
PROGRESS contains only present state and actionable work, and operational skills
link to those owners instead of repeating their content. Machine-enforced size
budgets prevent the former append-only growth pattern from returning. A new
\texttt{generate-report} skill now creates reproducible LaTeX and PDF reports
for Andrew at every wrap-up.

\section{Findings and results}
Before this change, the three primary files totaled 2,375 lines: 115 in
\texttt{CLAUDE.md}, 1,262 in the development plan, and 998 in PROGRESS. They
repeated requirements, benchmark protocol, current measurements, and historical
diagnoses. Several repetitions had diverged: N4 was omitted from one
non-negotiable list; the plan-revision workflow pointed to a history location
that had already moved; \texttt{SP-CHRONO-1} was both complete and open; and
stale T3 numbers appeared beside a statement that those numbers were void.

The same files are now 53, 342, and 167 lines respectively. The reduction keeps
the complete rule catalog, API, gates, open decisions, current phase, resume
state, and actionable parity backlog. Historical PROGRESS narrative was frozen
in \texttt{docs/log/2026-W32.md}. Remaining work identifiers include enough
context to act immediately and point to the self-contained parity ledger for
reference details and evidence.

\section{Evidence}
The new \texttt{tools/docs/check-doc-size.py} enforces budgets for living docs
and skills, verifies N1--N4 coverage, checks plan/PROGRESS revision agreement,
detects contradictory checklist states, and caps recent changes at five lines.
All documentation checks pass: 121 section references resolve, all 28 catalog
families have dispositions, and all 418 reference options remain accounted for.
The debug suite passes 59 of 59 tests with LeakSanitizer disabled because the
cluster session runs under ptrace; ASan and UBSan remain enabled.

The report toolchain uses the cluster-provided \texttt{amh-tex/2023} module.
Its build script compiles in a temporary directory, validates the PDF header,
and copies only the final PDF beside its LaTeX source.

\section{Risks and decisions}
Over-compression can make an identifier technically traceable but operationally
opaque. The chosen boundary is therefore: PROGRESS states the defect and next
action in one or two lines, while \texttt{docs/implementation-notes.md} owns the
reference mechanism, current implementation, evidence, and closer. Git history
is not required for ordinary resumption.

Andrew-facing reports add two frozen files per session by explicit decision.
They are excluded from living-document budgets and must synthesize rather than
copy logs. Missing LaTeX or a compilation failure blocks wrap-up instead of
leaving a source-only deliverable.

\section{Next actions}
The development resume point is unchanged by this documentation work: recover
the seven missing tasks from the fixed-binary Tier C ablation window, aggregate
the completed window, and use its source counters to attribute the three recent
ports. Thereafter, close the remaining N4 parity rows before the Tier D campaign
and T3 evaluation.

\end{document}
